\documentclass[11pt]{article}

\usepackage[T1]{fontenc}
\usepackage[utf8]{inputenc}
\usepackage[french]{babel}
\usepackage[a4paper,margin=2cm]{geometry}
\usepackage{amsmath,amssymb}
\usepackage{booktabs}
\usepackage[colorlinks=true,linkcolor=blue,urlcolor=blue]{hyperref}

\setlength{\parskip}{3pt}

\title{5OD21 -- Optimisation Discrète\\Projet : Course d'avions}
\author{Mohamed Benloughmari \and Fatimetou Narjah}
\date{}

\begin{document}
\maketitle
\vspace{-1cm}

\section{Problème et graphe}

On dispose de $n$ aérodromes $V=\{1,\dots,n\}$, d'un départ $d$, d'une arrivée $f$, de
$K$ régions (on note $V_k$ les aérodromes de la région $k$), d'un nombre minimal
d'aérodromes $A_{min}$ et d'un rayon d'action $R$. On cherche un chemin élémentaire de
$d$ à $f$, de longueur totale minimale, qui visite au moins $A_{min}$ aérodromes
(départ et arrivée compris), passe par chaque région et n'effectue que des sauts de
longueur $d_{ij}\le R$. C'est une variante du TSP orienté où l'on choisit les sommets
visités et où l'on cherche un chemin $d\to f$ plutôt qu'un circuit.

On travaille sur le graphe orienté $G=(V,A)$ avec
\[
A=\{(i,j)\in V^2 : i\neq j,\ d_{ij}\le R,\ j\neq d,\ i\neq f\}.
\]
La contrainte de rayon est ainsi prise en compte directement dans le graphe, et aucun
arc n'entre en $d$ ni ne sort de $f$. On note $\delta^+(k)$ (resp. $\delta^-(k)$) les arcs
sortant de (resp. entrant dans) $k$.

\section{Formulations}

\paragraph{Variables et contraintes communes.} $x_{ij}\in\{0,1\}$ vaut 1 si l'avion vole
directement de $i$ à $j$ ; $y_i\in\{0,1\}$ vaut 1 si l'aérodrome $i$ est visité.
\begin{align}
\min\ & \sum_{(i,j)\in A} d_{ij}\,x_{ij} \nonumber\\
\text{s.c.}\ & \textstyle\sum_{\delta^+(d)} x_{dj}=1,\qquad \sum_{\delta^-(f)} x_{if}=1,\qquad y_d=y_f=1 \label{c:df}\\
& \textstyle\sum_{\delta^-(k)} x_{ik}=y_k,\qquad \sum_{\delta^+(k)} x_{kj}=y_k \qquad \forall k\in V\setminus\{d,f\} \label{c:flot}\\
& \textstyle\sum_{i\in V} y_i\ge A_{min} \label{c:amin}\\
& \textstyle\sum_{i\in V_k} y_i\ge 1 \qquad \forall k\in\{1,\dots,K\} \label{c:reg}
\end{align}
Les contraintes \eqref{c:df}--\eqref{c:flot} imposent que les arcs choisis forment un
chemin de $d$ à $f$, auquel peuvent s'ajouter des circuits disjoints (sous-tours).
Ces sous-tours permettent de « visiter » des aérodromes à bas coût pour satisfaire
\eqref{c:amin}--\eqref{c:reg}, il faut donc les éliminer. Les deux formulations
diffèrent uniquement sur ce point.

\paragraph{Formulation 1 : MTZ (nombre polynomial de contraintes).} On ajoute une variable
continue $u_i\in[1,n]$ représentant la position de $i$ dans le chemin :
\begin{equation}
u_d=1,\qquad u_j\ge u_i+1-n\,(1-x_{ij})\quad \forall (i,j)\in A. \label{c:mtz}
\end{equation}
Si $x_{ij}=1$ on obtient $u_j\ge u_i+1$ ; en sommant le long d'un circuit
$i_1\to\dots\to i_p\to i_1$ on aurait $0\ge p$, ce qui est impossible. Si $x_{ij}=0$, la
contrainte est toujours satisfaite. Le modèle a $|A|+2n$ variables et $O(|A|)$
contraintes. En revanche, sa relaxation continue est faible : pour $x_{ij}$
fractionnaire, le terme $n(1-x_{ij})$ rend \eqref{c:mtz} presque inactive.

\paragraph{Formulation 2 : coupes de connexité (nombre exponentiel de contraintes).}
On remplace \eqref{c:mtz} par
\begin{equation}
\sum_{\substack{(i,j)\in A\\ i\notin S,\ j\in S}} x_{ij}\ \ge\ y_m
\qquad \forall S\subseteq V\setminus\{d\},\ \forall m\in S. \label{c:cut}
\end{equation}
Si $m$ est visité, tout ensemble $S$ contenant $m$ mais pas $d$ doit recevoir un arc :
$m$ est donc atteignable depuis $d$. Un sous-tour $C$ viole \eqref{c:cut} pour $S=C$.
C'est l'analogue des contraintes d'élimination de sous-tours du TSP, avec $y_m$ au second
membre car les sommets ne sont pas tous visités. Grâce à \eqref{c:flot}, elle est
équivalente à la forme GSEC
$\sum_{i,j\in S} x_{ij}\le \sum_{i\in S}y_i-y_m$. Le nombre de contraintes est de l'ordre
de $n\,2^{n-1}$ : on les ajoute uniquement lorsqu'elles sont violées.

\paragraph{Problème de séparation.} Étant donné $(x^*,y^*)$ satisfaisant
\eqref{c:df}--\eqref{c:reg}, trouver $S\subseteq V\setminus\{d\}$ et $m\in S$ tels que
$\sum_{i\notin S,\,j\in S}x^*_{ij}<y^*_m$, ou montrer qu'il n'en existe pas.
\begin{itemize}
\item \emph{Cas fractionnaire.} En prenant $x^*$ comme capacités, on a pour $m$ fixé
$\min_{S\ni m,\ d\notin S}\sum_{i\notin S,\,j\in S}x^*_{ij}$ égal à la coupe minimale
$d\to m$, donc au flot maximal de $d$ à $m$ (théorème flot max / coupe min). Pour chaque
$m$ tel que $y^*_m>0$, on calcule ce flot (Edmonds--Karp). S'il est strictement inférieur
à $y^*_m$, l'ensemble $S$ des sommets non atteignables depuis $d$ dans le graphe résiduel
fournit une coupe violée. La séparation est donc exacte et polynomiale ($n$ calculs de flot).
\item \emph{Cas entier.} Les arcs choisis forment un chemin et des circuits. On suit le
chemin depuis $d$ ; chaque circuit $C$ non atteint donne les coupes violées $(S=C,m)$
pour tout $m\in C$.
\end{itemize}

\section{Mise en œuvre}

Les modèles sont écrits en Julia avec JuMP et résolus avec SCIP, avec une limite de
600\,s par résolution. Pour la formulation 2, SCIP.jl ne prend pas en charge les
\emph{lazy callbacks} de JuMP (le code les utilise si l'on choisit Gurobi). Nous
utilisons donc une génération de contraintes itérative : on résout le PLNE, on sépare les
sous-tours de la solution entière, on ajoute les coupes, puis on recommence jusqu'à
obtenir une solution sans sous-tour. Celle-ci est optimale, puisqu'elle est optimale
pour un relâché du modèle complet et réalisable pour ce dernier. Le temps et le nombre
de nœuds indiqués sont cumulés sur les itérations.

La borne de relaxation continue à la racine ($z_{LP}$) est obtenue directement pour MTZ.
Pour la formulation 2, on applique un algorithme de plans coupants sur le PL relâché,
avec la séparation fractionnaire exacte, jusqu'à ce qu'aucune coupe ne soit violée.
Les coupes ainsi trouvées sont conservées dans le PLNE de départ, et le temps de cette
phase est inclus dans le temps total et dans la limite de 600\,s. Chaque solution est
vérifiée indépendamment (chemin $d\to f$ élémentaire, au moins $A_{min}$ sommets, sauts
$\le R$, toutes les régions visitées). Sur l'exemple du sujet (6 aérodromes), on retrouve
l'optimum 12 avec le parcours $1\to3\to2\to5$. Pour \texttt{aerodrome\_90\_1}, qui
déclare $K=6$ mais attribue la région 7 à 5 aérodromes, nous respectons $K$ et
considérons ces aérodromes comme hors région.

\section{Résultats}

Le tableau~\ref{tab:res} donne, pour chaque instance, la meilleure valeur connue $z^*$
(optimale sauf mention), la borne de relaxation continue à la racine $z_{LP}$, le gap
racine $(z^*-z_{LP})/z^*$, le temps total et le nombre de nœuds du Branch-and-Cut.
Pour la formulation 2, le nombre de nœuds est cumulé sur les itérations.

\begin{table}[h]
\centering\small
\setlength{\tabcolsep}{4.5pt}
\begin{tabular}{lrr|rrrr|rrrr}
\toprule
& & & \multicolumn{4}{c|}{MTZ} & \multicolumn{4}{c}{Coupes de connexité}\\
Instance & $|A|$ & $z^*$ & $z_{LP}$ & gap\,\% & temps (s) & nœuds & $z_{LP}$ & gap\,\% & temps (s) & nœuds\\
\midrule
0      &   18 &  12 &   8.0 & 33.3 &   0.0 &      1 &  10.0 & 16.7 &   0.3 & 1\\
20\_1  &  288 &  73 &  65.2 & 10.8 &   0.0 &      1 &  71.0 &  2.7 &   0.1 & 1\\
20\_2  &  286 &  91 &  63.0 & 30.8 &   0.7 &     67 &  74.5 & 18.1 &   0.0 & 1\\
20\_3  &  113 &  82 &  64.0 & 22.0 &   0.0 &      1 &  72.5 & 11.6 &   0.0 & 1\\
20\_4  &  226 &  96 &  78.0 & 18.8 &   0.1 &     10 &  86.5 &  9.9 &   0.0 & 1\\
30\_1  &  156 & 115 &  84.1 & 26.9 &   0.6 &    623 & 110.5 &  3.9 &   0.9 & 2\\
40\_1  &  417 & 124 & 114.2 &  7.9 &   0.0 &      1 & 124.0 &  0.0 &   0.6 & 1\\
50\_1  &  441 & 148 & 137.0 &  7.4 &   0.3 &      1 & 148.0 &  0.0 &   1.6 & 1\\
70\_1  &  757 & 288 & 231.1 & 19.8 &   7.7 &   1754 & 278.8 &  3.2 &  10.3 & 1\\
80\_1  &  608 & 385 & 221.3 & 42.5 & $>$600 & 491822 & 385.0 &  0.0 &  27.4 & 1\\
90\_1  & 1287 & 605 & 409.8 & 32.3 & $>$600 & 258278 & 604.3 &  0.1 &  62.0 & 2\\
100\_1 &  957 & 665 & 429.5 & 35.4 & $>$600 & 186610 & 658.7 &  0.9 &  54.4 & 2\\
100\_2 & 1205 & 580 & 410.9 & 29.1 & $>$600 & 208444 & 577.9 &  0.4 &  61.0 & 2\\
120\_1 & 5473 & 481$^\dagger$ & 248.6 & 48.3 & $>$600 & 67337 & 410.0 & 14.8 & $>$600 & 0\\
150\_1 & 1098 & 448 & 352.2 & 21.4 & $>$600 & 243733 & 446.9 &  0.2 & 294.2 & 1\\
\bottomrule
\end{tabular}
\caption{Résultats avec SCIP (limite 600\,s). $^\dagger$ meilleure solution trouvée,
optimalité non prouvée : on sait seulement que $z^*\in[410.0,\ 481]$.}
\label{tab:res}
\end{table}

Lorsque MTZ atteint la limite de temps, la meilleure solution trouvée et la borne
inférieure finale sont, pour (80\_1, 90\_1, 100\_1, 100\_2, 120\_1, 150\_1) :
(385 ; 367.8), (689 ; 545.0), (690 ; 639.1), (611 ; 559.3), (481 ; 388.7) et (456 ; 437.1).

\section{Analyse}

\paragraph{Qualité des relaxations.} Sur toutes les instances, la borne racine de la
formulation par coupes est meilleure que celle de MTZ. Le gap moyen passe d'environ
26\,\% (MTZ) à environ 5.5\,\% (coupes), et il est inférieur à 1\,\% sur toutes les
instances résolues avec $n\ge 80$.
Le gap de MTZ, au contraire, ne diminue pas quand $n$ augmente : il dépasse 30\,\% sur la
plupart des grandes instances. C'est le résultat attendu : avec $x_{ij}$ fractionnaire, la
contrainte $u_j\ge u_i+1-n(1-x_{ij})$ est presque toujours satisfaite, alors que les
coupes de connexité imposent, même en fractionnaire, qu'un flot $y_m$ puisse être
envoyé de $d$ vers chaque sommet $m$. On retrouve
la hiérarchie connue pour le TSP : la relaxation de la formulation exponentielle est plus
forte que celle de MTZ.

\paragraph{Effort de résolution.} Pour $n\le 70$, les deux formulations résolvent les
instances en quelques secondes. MTZ peut toutefois demander beaucoup de nœuds (1754 pour
70\_1), alors que la formulation par coupes termine presque toujours à la racine (1 ou 2
nœuds). À partir de $n=80$, MTZ ne prouve plus l'optimalité en 600\,s malgré plusieurs
centaines de milliers de nœuds. Sa meilleure solution peut rester loin de l'optimum :
689 contre 605 pour 90\_1, soit 14\,\% de plus. La formulation par coupes résout ces
instances en 30\,s à 5\,min. Le temps est alors presque entièrement consacré aux
plans coupants de la racine (environ $n$ calculs de flot maximal et un PL par itération) :
c'est le prix d'une relaxation forte.

\paragraph{Limites.} Le coût de la racine croît avec le nombre d'arcs. Sur 120\_1
($R=90$, donc 5473 arcs, graphe presque complet), les plans coupants épuisent à eux seuls
la limite de temps, et aucune solution entière n'est obtenue par cette formulation. Sa
borne (410.0) reste néanmoins meilleure que la borne finale de MTZ après 600\,s (388.7).
Deux améliorations simples seraient de limiter le temps passé à la racine, ou d'utiliser
un solveur qui accepte les \emph{lazy callbacks} (Gurobi) pour séparer les coupes au
cours du Branch-and-Cut plutôt que par résolutions successives. Le code prévoit déjà
cette seconde option.

\paragraph{Conclusion.} MTZ est simple à écrire et suffit pour les petites instances,
mais sa relaxation faible le rend inutilisable au-delà de 80 aérodromes. La formulation
exponentielle, associée à une séparation exacte par flot maximal, donne une borne
quasiment égale à l'optimum et résout 14 des 15 instances à l'optimalité, contre 9 pour
MTZ.

\end{document}
